People with amyotrophic lateral sclerosis progressively lose speech and movement but usually retain control of their eyes, yet commercial gaze-based communication devices are beyond the reach of most Brazilian families. This work presents a gaze-controlled communication app for ordinary Android tablets that relies only on the device's front camera and runs entirely offline. Eye tracking combines eye-corner detection with ML Kit, pupil localization with a YOLO26n network, a per-user polynomial regression calibration and a Kalman filter, and maps gaze direction to navigation and selection commands on an on-screen keyboard. To reduce typing effort, a trigram language model with Stupid Backoff, trained on augmentative and alternative communication phrases and tolerant to typing errors, suggests words, and the composed message is spoken by an on-device Portuguese text-to-speech model. A metrics module logs every session and computes typing speed, effort and the contribution of word suggestions. The system was evaluated in three usage tests performed by one of the authors on a Samsung Galaxy Tab S7, each repeated five times with every command issued by gaze. Performance improved with practice, and speed came to depend mainly on the number of commands rather than on the time per command. With word suggestions, command savings reached 79.4\% and text entry rate reached about 3 words per minute, comparable to other low-cost camera-based systems. Future work includes tests with people with amyotrophic lateral sclerosis, spoken feedback for the selected key and context-aware word suggestions.
